\documentclass[11pt,a4paper]{article}
\usepackage[T1]{fontenc}
\usepackage{lmodern}
\usepackage[margin=28mm]{geometry}
\usepackage{amsmath,amssymb,amsthm,mathtools}
\usepackage{booktabs,microtype}
\usepackage[hidelinks]{hyperref}
\hypersetup{pdftitle={Quadratic standard monomials for fixed involution classes},
 pdfauthor={Xiaojian Zeng}}
\newtheorem{theorem}{Theorem}
\newtheorem{lemma}[theorem]{Lemma}
\newtheorem{corollary}[theorem]{Corollary}
\newcommand{\kk}{\mathbb{K}}
\newcommand{\QQ}{\mathbb{Q}}
\newcommand{\CC}{\mathbb{C}}
\newcommand{\Mn}{\mathcal{M}_{n,a}}
\newcommand{\HH}{\mathcal{H}_n}
\newcommand{\TT}{\mathcal{T}_n}
\newcommand{\BB}{\mathcal{B}_{n,2}}
\DeclareMathOperator{\gr}{gr}
\DeclareMathOperator{\RSK}{RSK}
\DeclareMathOperator{\Hilb}{Hilb}
\DeclareMathOperator{\rank}{rank}
\DeclareMathOperator{\score}{s}
\title{Quadratic standard monomials for fixed involution classes}
\author{Xiaojian Zeng\thanks{PBC School of Finance, Tsinghua University,
100190, China. Email: \href{mailto:xiaojianzeng@foxmail.com}{xiaojianzeng@foxmail.com}.}}
\date{}
\begin{document}
\maketitle
\begin{abstract}
Let $\mathcal{M}_{n,a}$ be the permutation matrices of involutions with
$a$ fixed points. For $a\geq2$ and $(n-a)/2\geq2$, we prove that Zhu's
Pieri--RSK candidate is a basis of its orbit-harmonics quotient in degrees
$0$, $1$, and $2$. The quadratic basis consists of all two-edge matching
monomials except an explicit set of $\binom n2$ monomials. These exclusions
form a connected spanning subgraph of $KG(n,2)$ with one odd cycle.
Their unique straightening relations strictly decrease the variable weights
$w(y_{ij})=n^2-ij$, so the basis is standard for every monomial order
refining these weights in the off-diagonal presentation. In particular,
this gives a complete standard monomial basis for the class with exactly
two transpositions, for every $n\geq6$. The general monomial-basis problem
for arbitrary fixed involution classes remains unresolved by this result.
\end{abstract}

\section{The quotient and the result}

For integers $n>0$, $0\leq a\leq n$, and $a\equiv n\pmod2$, put
$k=(n-a)/2$. Write $[n]=\{1,\ldots,n\}$ and let
\[
 \Mn=\{P_w:w\in\mathfrak S_n,\ w^2=1,\
                |\{i:w(i)=i\}|=a\},\qquad
 (P_w)_{ij}=\mathbf1_{w(i)=j}.
\]
Work over $\kk=\QQ$ or $\CC$. If $I(\Mn)$ is the vanishing ideal in
$A=\kk[x_{ij}:1\leq i,j\leq n]$, define
\[
 \gr I(\Mn)=\langle\tau(f):0\ne f\in I(\Mn)\rangle,
 \qquad R_{n,a}=A/\gr I(\Mn),
\]
where $\tau(f)$ is the highest nonzero homogeneous part of $f$.
Thus the ideal is formed from all vanishing polynomials, and $R_{n,a}$
is graded by total degree. The known presentation recalled below is
defined over $\QQ$, and its complex scalar extension gives the same
construction over $\CC$.

Liu, Ma, Rhoades, and Zhu established a matching-monomial basis for the
locus of all involutions and determined the graded characters for a fixed
class \cite[Theorems 3.3 and 5.21]{LMRZ}.
Zhu's positive character formula leads to a particular Pieri--RSK
candidate for each graded component \cite[Proposition 17 and the
discussion following Problem 29]{ZhuPositive}. Problem 29 asks for an
explicit monomial basis of $R_{n,a}$ for general $n,a$. We prove the
candidate in a restricted range of degrees and obtain a complete basis
for the infinite family $a=n-4$.

Let $E=\binom{[n]}2$, and write $ij$ for the vertex $\{i,j\}\in E$,
always with $i<j$. Use off-diagonal variables $y_{ij}$, and let
\[
 S=\kk[y_{ij}:ij\in E],\qquad
 J_{n,a}=\ker\bigl(S\longrightarrow R_{n,a},\ y_{ij}\longmapsto[x_{ij}]\bigr).
\]
This map is onto by the linear relations below. Denote by $\HH$ the
set of squarefree two-edge matching monomials $y_e y_f$, $e\cap f=\varnothing$.
For $n\geq6$ define the following disjoint union:
\begin{equation}\label{eq:T}
\begin{aligned}
 \TT={}&\{y_{12}y_{ij}:3\leq i<j\leq n\}\\
 &\mathbin{\cup}\{y_{13}y_{2j}:4\leq j\leq n\}\\
 &\mathbin{\cup}\{y_{1j}y_{23}:4\leq j\leq n\}\\
 &\mathbin{\cup}\{y_{15}y_{24},\ y_{15}y_{34},\ y_{25}y_{34}\}.
\end{aligned}
\end{equation}
Set $\BB=\HH\setminus\TT$ and
$\score(y_{ij}y_{rs})=ij+rs$.

\begin{theorem}\label{thm:main}
Suppose $a\geq2$ and $k=(n-a)/2\geq2$, in addition to the admissibility
conditions above. The literal Pieri--RSK candidate of
\cite[Problem 29]{ZhuPositive}, recalled in Section~\ref{sec:tableaux},
has the following degree-$0$, degree-$1$, and degree-$2$ sets:
\begin{equation}\label{eq:bases}
 \{1\},\qquad \{y_{ij}:ij\in E,\ ij\ne12\},\qquad \BB.
\end{equation}
Their residue classes are bases of the corresponding components of
$R_{n,a}$. They are also exactly the standard monomial sets in those
degrees for $S/J_{n,a}$, for every monomial order refining
\begin{equation}\label{eq:weight}
 w(y_{ij})=n^2-ij,
\end{equation}
where larger weight means larger monomial.
For every $t\in\TT$ there is a unique equality in $R_{n,a}$
\begin{equation}\label{eq:straighten}
 t=\sum_{h\in\BB}c_{t,h}h,
 \qquad c_{t,h}\in\QQ,
 \qquad c_{t,h}\ne0\ \Longrightarrow\ \score(t)<\score(h).
\end{equation}
\end{theorem}

Here a monomial is standard if it does not belong to the initial ideal.
The order assertion concerns $S$, with its off-diagonal variables; it
does not prescribe an order on all $n^2$ original variables. The weights
are positive, since $n^2-ij\geq n$ for $i<j\leq n$.

\begin{corollary}\label{cor:two}
For every $n\geq6$, the union of the three sets in \eqref{eq:bases}
is a complete standard monomial basis of $R_{n,n-4}$ for every order in
Theorem~\ref{thm:main}. Its Hilbert series is
\begin{equation}\label{eq:hilbert}
 \Hilb(R_{n,n-4};q)
 =1+\left(\binom n2-1\right)q
    +\left(3\binom n4-\binom n2\right)q^2.
\end{equation}
\end{corollary}

The ideal presentation and graded dimensions are prior results
\cite{LMRZ,ZhuPositive,ZhuRook}. The relation between degree-one
Lefschetz multiplication and signless graph incidence, including its
rank criterion, is also known \cite[Proposition 3.1, Lemma 3.2 and
Theorem 3.3]{DaoNair}. Our contribution is the identification of the
specific tableau exclusions with \eqref{eq:T}, together with the strict
comparison for \eqref{eq:weight}. No equivalent of this precise combination
was found in the primary sources inspected through 14 September 2026;
this is a bounded literature assessment, not a priority claim.

\section{Quadratic relations and an odd cycle}\label{sec:graph}

We first recall the algebraic input. The presentation in
\cite[Section 2.5]{ZhuPositive}, proved in
\cite[Appendix A, Proposition A.7]{ZhuRook}, generates $\gr I(\Mn)$ by
the row sums, column sums, differences $x_{ij}-x_{ji}$, the trace
$\sum_i x_{ii}$, all quadratic products within a single row or column
(including squares), the products $\prod_{i\in U}x_{ii}$ with $|U|=a+1$,
and all monomials of degree greater than $k$. In the last family the
matching monomials suffice, as in the presentation of \cite{ZhuRook}.

Consider the squarefree matching algebra and its linear form
\begin{equation}\label{eq:B}
 B=S/(y_e^2,\ y_e y_f:e\ne f,\ e\cap f\ne\varnothing),
 \qquad L=\sum_{e\in E}y_e.
\end{equation}
The symmetry and row-sum relations eliminate the remaining coordinates:
\begin{equation}\label{eq:elimination}
 x_{ji}=y_{ij}\ (i<j),\qquad
 x_{ii}=-\sum_{j\ne i}y_{\min(i,j),\max(i,j)}.
\end{equation}
Column sums then vanish as well. Every row or column quadratic becomes
zero in $B$: two off-diagonal factors share an endpoint, and terms
involving $x_{ii}$ expand into such products, including squares.
Conversely, each defining quadratic of $B$ comes from a row or column
product using symmetry. The trace becomes $-2L$.
The remaining generators have degrees $a+1\geq3$ or greater than
$k\geq2$. All substitutions are linear and homogeneous, so these
generators have no consequences in smaller degrees. Hence
\begin{equation}\label{eq:low}
 (R_{n,a})_0=\kk,\qquad
 (R_{n,a})_1=B_1/\kk L,\qquad
 (R_{n,a})_2=B_2/LB_1.
\end{equation}

Let $G=KG(n,2)$ be the graph whose vertices are $E$ and whose edges join
disjoint pairs. Its edges index $\HH$, the monomial basis of $B_2$.
Let $D$ be its signless vertex--edge incidence matrix: the column of an
edge $ef$ is $\mathbf e_e+\mathbf e_f$. The relation
\[
 y_e L=\sum_{f\cap e=\varnothing}y_e y_f
\]
shows that the rows of $D$ span precisely $LB_1$ in these coordinates.
This is the incidence interpretation of \cite[Lemma 3.2]{DaoNair}.

\begin{lemma}\label{lem:incidence}
The submatrix $D_T$ formed by the columns in $\TT$ is square and
invertible over $\QQ$, with $|\det D_T|=2$. Consequently $\BB$ is a
basis of $(R_{n,a})_2$, and its straightening coefficients are the entries
of $-D_T^{-1}D_{\BB}$.
\end{lemma}

\begin{proof}
The first three families of \eqref{eq:T} form three vertex-disjoint
stars covering $E$: their centers are $12,13,23$, respectively.
The last three edges connect the stars. The resulting spanning graph
is connected and has
\[
 \binom{n-2}{2}+2(n-3)+3=\binom n2
\]
edges, as many as vertices. Its unique cycle is
\begin{equation}\label{eq:cycle}
 13\;-\;24\;-\;15\;-\;34\;-\;25\;-\;13.
\end{equation}
If $z^\mathsf{t}D_T=0$, the entries of $z$ have opposite values at
adjacent vertices. The odd cycle forces them to be zero there, and
connectedness forces $z=0$. This proves invertibility. Removing leaf
rows and their incident columns reduces the determinant, up to sign,
to that of the five-cycle, which has absolute value $2$.

Writing the relations as a vector equality in the quotient gives
\[
 D_T(t)_{t\in\TT}+D_{\BB}(h)_{h\in\BB}=0.
\]
The relation space projects isomorphically onto the $\TT$ coordinates,
so the complementary monomials are independent and span the quotient
in \eqref{eq:low}. Solving this system gives the asserted coefficients.
\end{proof}

\section{Strict comparison of weights}\label{sec:weights}

For a graph vertex $v=ij$ put $p(v)=ij$. Thus
$\score(uv)=p(u)+p(v)$ for a graph edge. The following column statement
will determine every order refining \eqref{eq:weight}.

\begin{lemma}\label{lem:score}
For each $h=uv\in\BB$, the vector $\mathbf e_u+\mathbf e_v$ is in the
span of the incidence columns of those $t\in\TT$ with
$\score(t)<\score(h)$.
\end{lemma}

\begin{proof}
The alternating sum of incidence columns along an odd walk from $u$ to
$v$ is $\mathbf e_u+\mathbf e_v$. This identity permits repeated edges
or vertices. Also, the columns of a connected graph containing an odd
cycle span its entire vertex coordinate space, by the alternating-value
argument in Lemma~\ref{lem:incidence}.

Every edge incident with vertex $12$ belongs to $\TT$. The other edges
$h$ fall into the following cases.

\smallskip\noindent
\emph{Case 1: $h=(1r)(2s)$.}
If $r=3$ or $s=3$, then $h\in\TT$. Otherwise $r,s\geq4$, $r\ne s$,
and the odd walk
\[
 1r\;-\;23\;-\;15\;-\;24\;-\;13\;-\;2s
\]
uses only edges of $\TT$. Its edge scores are
$r+6,11,13,11,2s+3$. The first and last are strictly less than $r+2s$.
The minimum possible $r+2s$ is $13$, attained only at $(r,s)=(5,4)$;
that edge is in $\TT$. Thus all five scores are strictly smaller than
$\score(h)$ for the edges under consideration.

\smallskip\noindent
\emph{Case 2: $h=(1r)(st)$, with $3\leq s<t$ and $2$ absent.}
When $r\geq4$, use the odd walk
\[
 1r\;-\;23\;-\;15\;-\;34\;-\;12\;-\;st.
\]
Its scores are $r+6,11,17,14,st+2$. Since $st\geq12$ and $r\geq4$,
the endpoint comparisons are strict. Disjointness gives $r+st\geq19$
if $r=4$. For $r=5$, equality at the minimum $17$ occurs only for
$15$--$34$, which belongs to $\TT$. If $r\geq6$, then $r+st\geq18$.
All comparisons are therefore strict. The exception $r=3$ has score
at least $3+4\cdot5=23$ and is included in the next argument.

\smallskip\noindent
\emph{Case 3: the remaining edges.}
Apart from that exception, they have type $(2r)(st)$ with $s,t\geq3$,
or contain neither $1$ nor $2$. In the first type, the lower bounds
for $r=3,4,5,$ and $r\geq6$ are, respectively, $26,23,22,$ and $24$.
Equality at $22$ occurs only for $25$--$34$, an edge of $\TT$.
In the second type the minimum is $3\cdot6+4\cdot5=38$.
Thus all remaining $h$ have score greater than $22$.

The restriction of $\TT$ to the ten vertices $ij\subseteq[5]$ is
connected, contains \eqref{eq:cycle}, and has maximum edge score $22$.
Each endpoint of $h$ outside this core is attached to the core by one
edge of $\TT$: attach $1r$ to $23$, $2r$ to $13$, and $ij$ with
$3\leq i<j$ to $12$. These chosen neighbors uniquely minimize $p$
among vertices disjoint from the endpoint: their products are $6,3,$
and $2$, respectively. Since $h\notin\TT$, the attachment is not $h$,
and its score is strictly smaller than $\score(h)$. The subgraph formed
by this core and the necessary attachments is connected, contains an odd
cycle, and contains both endpoints of $h$. Its columns therefore span
$\mathbf e_u+\mathbf e_v$ using only strictly smaller scores.
\end{proof}

By invertibility of $D_T$, the representation in Lemma~\ref{lem:score}
is unique. A nonzero entry at row $t$ and column $h$ of
$D_T^{-1}D_{\BB}$ consequently implies $\score(t)<\score(h)$.
This proves \eqref{eq:straighten}. Since the quadratic weight is
$2n^2-\score$, each straightening relation has the unique largest-weight
monomial $t$ as its leading monomial. Every nonmatching quadratic also
belongs individually to $J_{n,a}$.

These leading monomials leave just $\BB$ in degree two, and
Lemma~\ref{lem:incidence} gives
$\dim(R_{n,a})_2=|\BB|$. There can be no further degree-two initial
monomial: in each fixed degree, row reduction of the homogeneous ideal
with columns ordered by the chosen monomial order leaves exactly the
dimension of the quotient as nonpivot monomials. This argument applies
to every monomial order, without a degree-compatibility assumption.
In degree one the only relation is $L$, whose unique largest-weight term
is $y_{12}$. The degree-zero assertion is immediate. This proves the
algebraic and order parts of Theorem~\ref{thm:main}.

\section{The literal Pieri--RSK candidate}\label{sec:tableaux}

For precision, we recall the candidate with English Young diagrams and
ordinary row insertion. To insert a label, replace the leftmost larger
entry of the first row and insert the displaced entry into the next row;
append if no larger entry exists. The RSK matrix--tableau correspondence
uses this insertion \cite[Sections 2--4]{Knuth}.

For $0\leq d\leq k$, take a partition $\mu$ of $2d$ and a partition $\lambda$ of $n$
such that every row length of $\mu$ is even,
$\lambda_1\leq n-2d+a$, and $\lambda/\mu$ is a horizontal strip
(at most one strip box in each column). Starting at height zero, scan
the columns of $\lambda$ from left to right, adding $+1$ when a column
meets the strip and $-1$ otherwise. The strip is positive if all the
heights through column $\lambda_1$ are nonnegative
\cite[Definition 5 and equation (4)]{ZhuPositive}.

For each standard tableau $P$ of shape $\lambda$, remove the strip by
reverse row insertion, starting at its rightmost column. In each preceding
row, reverse bump through the rightmost smaller entry. If the remaining
tableau is $Q$, the candidate monomial is obtained from
\begin{equation}\label{eq:inverse}
 M=\RSK^{-1}(Q^{\mathsf t},Q^{\mathsf t}),\qquad
 m=\prod_{i<j:M_{ij}=1}y_{ij}.
\end{equation}
Unused labels index zero rows and columns of $M$. The even row lengths
of $Q$ make $M$ a symmetric partial permutation matrix with zero diagonal
\cite[Proposition 2]{ZhuPositive}. The transpose in \eqref{eq:inverse}
is part of the construction.

Equivalently, start with a matching, apply RSK to its partial permutation
matrix, transpose the common tableau, and insert the unused labels in
increasing order. The inserted boxes give $\lambda/\mu$. This is the
inverse of the preceding procedure by the tableau version of Pieri's
rule \cite[Proposition 3]{ZhuPositive}. It suffices to test the resulting
strip and first-row restriction.

\begin{lemma}\label{lem:tableaux}
For $a\geq2$ and $k\geq2$, the degree-two candidate is exactly $\BB$.
In degree one it excludes only $y_{12}$, and in degree zero it is $\{1\}$.
\end{lemma}

\begin{proof}
For a two-edge matching on support $s_1<s_2<s_3<s_4$, four-letter row
insertion and transposition give the following three starting tableaux:
\begin{equation}\label{eq:tableaux}
\begin{array}{c|c}
 \text{matching}&Q\\ \hline
 (s_1s_2)(s_3s_4)&
    \begin{matrix}s_1&s_2\\s_3&s_4\end{matrix}\\[6pt]
 (s_1s_3)(s_2s_4)&
    \begin{matrix}s_1&s_3\\s_2&s_4\end{matrix}\\[6pt]
 (s_1s_4)(s_2s_3)&
    \begin{matrix}s_1&s_2&s_3&s_4\end{matrix}
\end{array}
\end{equation}
For either tableau of shape $(2,2)$, positivity holds precisely when
the strip meets column one, or equivalently when insertion creates a
third row. Indeed, the first step must be up; the second step then
cannot pass below zero; and every subsequent column is a strip column.
In the first tableau a third row is created if and only if the smallest
unused label is less than $s_2$. Such an insertion bumps a first-row
entry smaller than a second-row entry and so reaches row three.
If there is no such unused label, every insertion appends to row one.
Failure is therefore equivalent to $s_2=2$. In the second tableau
the same argument uses $s_3$, and failure is equivalent to $s_3=3$.
These are exactly the first two families in \eqref{eq:T}.

For the one-row tableau, increasing insertion bumps only original support
labels, in increasing order, into row two. Let its final second-row length
be $b$. It is the maximum number of disjoint pairs consisting of an unused
label and a larger support label: leftmost-greater insertion is the greedy
matching, and exchanging its first chosen support label for any larger
choice cannot improve the number of pairs remaining. The same exchange
argument applies successively. There is no third row. Among the first
four columns, the strip occupies exactly columns $1,\ldots,b$ in row two.
The path has $b$ up steps followed by $4-b$ down steps, and positivity is
equivalent to $b\geq2$; all later columns are strip columns.

Two such pairs exist precisely when
\begin{equation}\label{eq:two-pairs}
 s_3\geq4\quad\hbox{and}\quad s_4\geq6.
\end{equation}
For necessity, order the two paired support labels: the smaller is at
most $s_3$ and the larger at most $s_4$. Thus at least one unused label
is below $s_3$, and at least two are below $s_4$. Their respective counts
are $s_3-3$ and $s_4-4$. Conversely, these counts permit pairing the
smallest unused label with $s_3$ and the second smallest with $s_4$.
Failure of \eqref{eq:two-pairs} means $s_3=3$ or $s_4\leq5$.
The first condition gives $y_{1j}y_{23}$, $j\geq4$.
After removing these cases, the second condition leaves the supports
$\{1,2,4,5\}$, $\{1,3,4,5\}$, and $\{2,3,4,5\}$.
Their nested pairings are exactly the last three monomials of
\eqref{eq:T}.

The extra first-row restriction is automatic here by
\cite[Corollary 24]{ZhuPositive}, since $d=2\leq a$.
For degree one the starting tableau is one row with the two matched
labels. Positivity requires at least one bump into row two; it fails
exactly for support $\{1,2\}$. The same corollary applies since $1\leq a$.
In degree zero the empty tableau, followed by insertion of all labels,
gives the single monomial $1$.
\end{proof}

Lemma~\ref{lem:tableaux} completes Theorem~\ref{thm:main}.
For $a=n-4$ and $n\geq6$, we have $a\geq2$ and $k=2$.
All higher degrees vanish by the cited ideal presentation, or
\cite[Lemma A.3]{ZhuRook}. Thus the three sets form the complete standard
basis. There are $3\binom n4$ two-edge matchings and $\binom n2$
excluded ones, proving Corollary~\ref{cor:two}. The total basis size is
$3\binom n4$, as expected for the class with two transpositions.

\section{Exact checks and scope}

The proof is uniform and does not depend on computation. Exact finite
checks provide additional tests of its conventions and identities.
One implementation checks every quadratic matching for $6\leq n\leq14$,
including the incidence inverse, determinant, strict score comparisons,
and agreement of forward and inverse tableau constructions.
A separately written implementation solves the incidence systems by leaf
removal and a five-cycle recurrence, storing twice the coefficients as
integers. For $n=6,7,8,9,10,12,15,20$, every incidence identity and every
nonzero coefficient comparison passes; the minimum positive score margin
is $1$. Its independent enumeration of standard tableaux checks the literal
inverse candidate in degrees one and two for $6\leq n\leq10$.
These are finite observations, not an extension of the theorem's range.

As a separate check on the quotient, let $F_{\leq d}$ be the span of
evaluations on $\Mn$ of all polynomials of degree at most $d$.
Then $(R_{n,a})_d\cong F_{\leq d}/F_{\leq d-1}$.
For candidate columns $C_d$, the relevant test is
\[
 \rank[F_{\leq d-1},C_d]-\dim F_{\leq d-1}
 =|C_d|=\dim F_{\leq d}-\dim F_{\leq d-1}.
\]
This tests independence modulo lower-degree functions.
On the point set, symmetry,
$x_{ii}=1-\sum_{j\ne i}x_{ij}$, idempotence of off-diagonal coordinates,
and vanishing of overlapping edge products reduce every polynomial to
matching monomials without increasing degree. Exact rational ranks,
computed with SymPy~1.14.0 \cite{SymPy}, give Table~\ref{tab:ranks}.
The points were enumerated directly from permutations. The $(6,2)$
calculation also includes every original-coordinate monomial through
degree two; the other cases use the preceding spanning reduction.

\begin{table}[htbp]
\centering
\begin{tabular}{ccrrr}
\toprule
$(n,a)$ & $|\Mn|$ & $\dim F_{\leq1}$ & $\dim F_{\leq2}$
 & $\dim(R_{n,a})_2$\\
\midrule
$(6,2)$ & 45  & 15 & 45  & 30\\
$(7,3)$ & 105 & 21 & 105 & 84\\
$(8,2)$ & 420 & 28 & 210 & 182\\
\bottomrule
\end{tabular}
\caption{Exact finite evaluation checks. In every row $\dim F_{\leq0}=1$;
adjoining the quadratic candidate to $F_{\leq1}$ gives rank
$\dim F_{\leq2}$.}\label{tab:ranks}
\end{table}

The incidence checks retain exact matrix identities, so their verification
does not rely on an isolated numerical rank. A related use of reconstructible
multiplication-matrix certificates for monomial algebras is reported by
Zeng \cite{ZengWhiskered} for whiskered graphs. This is a
methodological comparison; no result about whiskered graphs is needed
for the proof here. Reproduction commands and the precise finite coverage
are recorded in the accompanying verification documentation.

The restrictions $a\geq2$ and $k\geq2$ are essential to the quadratic
reduction used here. We make no assertion for $a<2$ or about standard
monomials in arbitrary higher degrees. In particular, a uniform proof of
the full Pieri--RSK candidate, or of the same weight rule in all degrees,
is not supplied. The original general construction problem remains
unresolved by this note.

\paragraph{Computational and AI assistance.}
OpenAI Codex assisted proof exploration, preparation of the verification
code, argument checking in separate model contexts, and manuscript drafting.
These model passes are not external peer review. All reported calculations
use exact integer or rational arithmetic, with no random or floating-point
rank tests. The uniform argument is given above; no formal proof-assistant
certificate is claimed.

\bibliographystyle{amsplain}
\bibliography{references}
\end{document}
